\documentclass[a4paper,oneside]{report}
\usepackage{scrextend}
\changefontsizes{13pt}
\usepackage[utf8]{vietnam}
\usepackage[top=2cm, bottom=2cm, left=3.5cm, right=2.5cm]{geometry}

\usepackage{graphicx} % Cho phép chèn hỉnh ảnh
\usepackage{fancybox} % Tạo khung box
\usepackage{indentfirst} % Thụt đầu dòng ở dòng đầu tiên trong đoạn
\usepackage{amsthm} % Cho phép thêm các môi trường định nghĩa
\usepackage{latexsym} % Các kí hiệu toán học
\usepackage{amsmath} % Hỗ trợ một số biểu thức toán học
\usepackage{amssymb} % Bổ sung thêm kí hiệu về toán học
\usepackage{amsbsy} % Hỗ trợ các kí hiệu in đậm
\usepackage{times} % Chọn font Time New Romans
\usepackage{array} % Tạo bảng array
\usepackage{enumitem} % Cho phép thay đổi kí hiệu của list
\usepackage{subfiles} % Chèn các file nhỏ, giúp chia các chapter ra nhiều file hơn
\usepackage{titlesec} % Giúp chỉnh sửa các tiêu đề, đề mục như chương, phần,..
\usepackage{titletoc}
\usepackage{chngcntr} % Dùng để thiết lập lại cách đánh số caption,..
\usepackage{multirow} % Merge cells
\usepackage{longtable} % Bảng dài, ngắt qua trang
\usepackage{fancyhdr} % Cho phép tùy biến header và footer
\usepackage{appendix}

\usepackage[font=small,labelfont=bf]{caption}

\usepackage{listings}
\usepackage{float}
\usepackage{subcaption}
\usepackage{xurl}
\usepackage{booktabs,tabularx}
\usepackage[table]{xcolor}
\usepackage{adjustbox} % thu bảng rộng về vừa khổ giấy
\setlength{\emergencystretch}{3em} % tên file, tên tín hiệu dạng \texttt không ngắt được

\usepackage[nonumberlist, nopostdot, nogroupskip, acronym]{glossaries}
\usepackage{glossary-superragged}
\setglossarystyle{superraggedheaderborder}
\usepackage{setspace}
\usepackage{parskip}

% package content table
\usepackage{tocbasic}

% ===================================================

\usepackage[backend=biber,style=ieee]{biblatex}

\renewcommand\appendixname{PHỤ LỤC}
\renewcommand\appendixpagename{PHỤ LỤC}
\renewcommand\appendixtocname{PHỤ LỤC}

\addbibresource{\subfix{Danh_sach_tai_lieu_tham_khao.bib}}

\input{\subfix{Tu_viet_tat}}
\input{\subfix{generated/numbers}}

% Màu ô trong giản đồ pipeline và khung code
\definecolor{cIF}{HTML}{DCEBFA}
\definecolor{cID}{HTML}{E4F3DC}
\definecolor{cEX}{HTML}{FCE9C8}
\definecolor{cME}{HTML}{F3DDF0}
\definecolor{cWB}{HTML}{E2E2E2}
\definecolor{cST}{HTML}{F8C9C4}
\definecolor{code}{HTML}{F6F6F6}
\lstset{basicstyle=\ttfamily\small,backgroundcolor=\color{code},frame=none,columns=fullflexible,
  keepspaces=true,showstringspaces=false,xleftmargin=0.5em,aboveskip=0.5em,belowskip=0.5em,
  breaklines=true,breakatwhitespace=true,basicstyle=\ttfamily\footnotesize}

\newcommand{\sig}[1]{\texttt{#1}} % tên khối, tín hiệu trong mạch
\newcommand{\file}[1]{\texttt{#1}} % tên file

% ===================================================

\fancypagestyle{plain}{%
\fancyhf{} % clear all header and footer fields
\fancyfoot[RO,RE]{\thepage} %RO=right odd, RE=right even
\renewcommand{\headrulewidth}{0pt}
\renewcommand{\footrulewidth}{0pt}}

\setlength{\headheight}{18pt} % fancyhdr yêu cầu tối thiểu ~17.98pt với cỡ chữ 13pt

\def \TITLE{BÁO CÁO BÀI TẬP LỚN KIẾN TRÚC MÁY TÍNH TIÊN TIẾN}
\def \AUTHOR{Quyền Đức Công, Bùi Phạm Sơn Hà, Nguyễn Đức Anh}

% ===================================================
\titleformat
    {\chapter} % command
    [hang] % shape
    {\centering\bfseries} % format
    {CHƯƠNG \thechapter.\ } % label
    {0pt} %sep
    {} % before
    [] % after
\titlespacing*{\chapter}{0pt}{-20pt}{20pt}

\titleformat
    {\section} % command
    [hang] % shape
    {\bfseries} % format
    {\thechapter.\arabic{section}\ \ \ \ } % label
    {0pt} %sep
    {} % before
    [] % after
\titlespacing{\section}{0pt}{\parskip}{0.5\parskip}

\titleformat
    {\subsection} % command
    [hang] % shape
    {\bfseries} % format
    {\thechapter.\arabic{section}.\arabic{subsection}\ \ \ \ } % label
    {0pt} %sep
    {} % before
    [] % after
\titlespacing{\subsection}{30pt}{\parskip}{0.5\parskip}

\renewcommand\thesubsubsection{\alph{subsubsection}}
\titleformat
    {\subsubsection} % command
    [hang] % shape
    {\bfseries} % format
    {\alph{subsubsection}, \ } % label
    {0pt} %sep
    {} % before
    [] % after
\titlespacing{\subsubsection}{50pt}{\parskip}{0.5\parskip}

% ===================================================
\usepackage{hyperref}
\hypersetup{pdfborder = {0 0 0}} %
\hypersetup{pdftitle={\TITLE},
	pdfauthor={\AUTHOR}}

\usepackage[all]{hypcap} % Cho phép tham chiếu chính xác đến hình ảnh và bảng biểu

% Ảnh mạch chụp từ Logisim nằm ở docs/img; bảng và biểu đồ sinh tự động nằm ở generated/
\graphicspath{{../docs/img/}{generated/}{../../docs/img/}{../generated/}}

\counterwithin{figure}{chapter} % Đánh số hình ảnh kèm theo chapter. Ví dụ: Hình 1.1, 1.2,..

\title{\bf \TITLE}
\author{\AUTHOR}

\setcounter{secnumdepth}{3} % Cho phép subsubsection trong report

\theoremstyle{definition}
\newtheorem{example}{Ví dụ}[chapter] % Định nghĩa môi trường ví dụ

\onehalfspacing
%Khoảng cách xuống dòng
\setlength{\parskip}{6pt}
%Lùi đầu dòng
\setlength{\parindent}{15pt}

% =========================== BODY ===============
\begin{document}
\subfile{Bia} % Phần bìa

% ===================================================
\newpage
\pagenumbering{roman}
\renewcommand*\contentsname{MỤC LỤC}

\titlecontents{chapter}
    [0.0cm]             % left margin
    {\bfseries\vspace{0.3cm}}                  % above code
    {{\bfseries{\scshape}
    CHƯƠNG \thecontentslabel.\ }}
    % numbered format
    {}         % unnumbered format
    {\titlerule*[0.3pc]{.}\contentspage}         % filler-page-format, e.g dots

\titlecontents{section}
    [0.0cm]             % left margin
    {\vspace{0.3cm}}                  % above code
    {\thecontentslabel \ } % numbered format
    {}         % unnumbered format
    {\titlerule*[0.3pc]{.}\contentspage}         % filler-page-format, e.g dots

\titlecontents{subsection}
    [1.0cm]             % left margin
    {\vspace{0.3cm}}                  % above code
    {\thecontentslabel \ } % numbered format
    {}         % unnumbered format
    {\titlerule*[0.3pc]{.}\contentspage}         % filler-page-format, e.g dots

 % Tạo mục lục tự động
\addtocontents{toc}{\protect\thispagestyle{empty}}
\tableofcontents
\thispagestyle{empty}
\clearpage

%Tạo danh mục hình vẽ.
\renewcommand{\listfigurename}{DANH MỤC HÌNH VẼ}
{\let\oldnumberline\numberline
\renewcommand{\numberline}{Hình~\oldnumberline}
\listoffigures}

%Tạo danh mục bảng biểu.
\renewcommand{\listtablename}{DANH MỤC BẢNG BIỂU}
{\let\oldnumberline\numberline
\renewcommand{\numberline}{Bảng~\oldnumberline}
\listoftables}

\glsaddall
\renewcommand*{\glossaryname}{Danh sách thuật ngữ}
\renewcommand*{\acronymname}{DANH MỤC THUẬT NGỮ VÀ TỪ VIẾT TẮT}
\renewcommand*{\entryname}{Thuật ngữ}
\renewcommand*{\descriptionname}{Ý nghĩa}
\printnoidxglossary[type=\acronymtype]

% ===================================================

\newpage
\pagenumbering{arabic}

\pagestyle{fancy}
\fancyhf{}
\fancyhead[RE, LO]{\small\leftmark}
\fancyfoot[RE, LO]{\thepage}

\chapter{GIỚI THIỆU}
\label{ch:intro}
\subfile{chapters/01_gioi_thieu}

\chapter{TẬP LỆNH RV32I VÀ KIẾN TRÚC TỔNG THỂ}
\label{ch:arch}
\subfile{chapters/02_kien_truc}

\chapter{THIẾT KẾ CÁC MODULE}
\label{ch:module}
\subfile{chapters/03_module}

\chapter{CPU ĐƠN CHU KỲ}
\label{ch:single}
\subfile{chapters/04_single}

\chapter{CPU PIPELINE 5 TẦNG VÀ XỬ LÝ HAZARD}
\label{ch:pipe}
\subfile{chapters/05_pipeline}

\chapter{KIỂM THỬ}
\label{ch:test}
\subfile{chapters/06_kiem_thu}

\chapter{KẾT QUẢ VÀ ĐÁNH GIÁ}
\label{ch:result}
\subfile{chapters/07_ket_qua}

\chapter{PHẦN MỞ RỘNG, HẠN CHẾ VÀ HƯỚNG PHÁT TRIỂN}
\label{ch:ext}
\subfile{chapters/08_mo_rong}

% ===================================================
\newpage
\renewcommand\bibname{TÀI LIỆU THAM KHẢO}
\phantomsection\addcontentsline{toc}{chapter}{TÀI LIỆU THAM KHẢO}
\printbibliography

\appendixpage
\appendices
\addappheadtotoc

\titleformat{\chapter}[hang]{\centering\bfseries}{ \thechapter.\ }{0pt}{}[]
\titlespacing*{\chapter}{0pt}{-20pt}{20pt}

\titlecontents{chapter}
    [0.0cm]             % left margin
    {\bfseries\vspace{0.3cm}}                  % above code
    {{\bfseries{\scshape} \thecontentslabel.\ }} % numbered format
    {}         % unnumbered format
    {\titlerule*[0.3pc]{.}\contentspage}         % filler-page-format, e.g dots
\chapter{HƯỚNG DẪN SỬ DỤNG PROJECT}
\label{app:guide}
\subfile{chapters/A_huong_dan}
\end{document}
